\documentclass[12pt,a4paper]{article}

\usepackage[utf8]{inputenc}
\usepackage[T1]{fontenc}
\usepackage[spanish,provide=*]{babel}
\usepackage{graphicx}
\usepackage{amsmath}
\usepackage{amssymb}
\usepackage{geometry}
\usepackage{array}
\usepackage{verbatim}

\geometry{margin=2.5cm}

\newcommand{\code}[1]{\texttt{#1}}

\title{\textbf{An\'alisis Comparativo de Algoritmos de B\'usqueda:\\ A* frente a Hill Climbing en un Problema de Navegaci\'on}}
\author{Carlos Manuel Pi\~nero Ram\'irez \\ Facultad de Tecnolog\'ias Libres \\ Universidad de las Ciencias Inform\'aticas}
\date{\today}

\begin{document}
\maketitle

\begin{abstract}
\noindent
Se contrast\'o emp\'iricamente una b\'usqueda heur\'istica global (A*) frente a una b\'usqueda
local (Steepest-Ascent Hill Climbing) sobre un plano discreto de $15\times15$ celdas con
obst\'aculos, cargado desde \code{dataset\_coordenadas.csv}. A* result\'o \'optimo con un
costo de $24.00$ usando la distancia eucl\'idea como heur\'istica admisible y consistente,
expandiendo $152$ nodos. Hill Climbing alcanz\'o el mismo \'optimo expandiendo solo $59$
nodos, pero su \'exito dependi\'o de la topolog\'ia: un barrido desde las $206$ celdas
libres confirm\'o la ausencia de \'optimos locales en este tablero, mientras que un
escenario adversario en forma de ``U'' lo dej\'o atrapado en un m\'aximo local. Se concluye
que la b\'usqueda local es eficiente pero no confiable, y que A* paga su optimalidad con un
mayor n\'umero de expansiones.
\end{abstract}

\section{Introducci\'on}
\label{sec:intro}

El problema de navegaci\'on en un espacio de estados consiste en hallar una secuencia de
acciones que conduzca desde un estado inicial hasta un estado objetivo evitando regiones
prohibidas. Este trabajo aborda dicho problema en un plano 2D discreto y compara dos
familias de algoritmos con garant\'ias te\'oricas muy distintas.

La \textbf{b\'usqueda heur\'istica global} A* \cite{russell2020} es un algoritmo de
b\'usqueda informada que combina el costo real acumulado $g(n)$ con una estimaci\'on
$h(n)$ del costo restante. Cuando $h(n)$ es \emph{admisible} (nunca sobreestima el costo
real) A* es \textbf{\'optimo}, y cuando es \emph{consistente} es adem\'as eficiente en
t\'erminos de reexpansiones. En cambio, la \textbf{b\'usqueda local} Hill Climbing
\cite{russell2020} solo examina el vecindario inmediato del estado actual y avanza hacia
el vecino con mejor evaluaci\'on. Es un m\'etodo voraz, de bajo costo por paso, pero
carente de garant\'ias de completitud y optimalidad: puede quedar atrapado en
\textbf{m\'aximos locales}, \textbf{mesetas} o \textbf{crestas}.

El objetivo es comprender de forma emp\'irica c\'omo la topolog\'ia del espacio de
b\'usqueda determina el desempe\~no de cada algoritmo, evidenciando los conceptos de
completitud, optimalidad y \'optimos locales.

\section{Metodolog\'ia}
\label{sec:metodologia}

\subsection{Definici\'on del espacio de estados}
\label{sec:espacio}

El espacio se construy\'o a partir del archivo \code{dataset\_coordenadas.csv}, que
contiene $225$ registros con columnas \code{x}, \code{y} y \code{tipo}. Cada registro
define una celda del plano con uno de cuatro tipos: \code{libre}, \code{inicio},
\code{objetivo} u \code{obstaculo}. El espacio resultante es una cuadr\'icula de
$15\times15$ celdas con $206$ celdas libres, $17$ obst\'aculos, el estado inicial en
$(1,1)$ y el estado objetivo en $(13,13)$.

Formalmente, el espacio de estados es la tupla
\[
M = (S,\, s_0,\, S_g,\, A,\, c),
\]
donde:
\begin{itemize}
    \item $S$ es el conjunto de celdas transitables (libres y objetivo);
    \item $s_0 = (1,1)$ es el estado inicial y $S_g = \{(13,13)\}$ el estado objetivo;
    \item $A = \{\textsc{Arriba}, \textsc{Abajo}, \textsc{Izquierda}, \textsc{Derecha}\}$
          define 4-conectividad, es decir, sin movimientos diagonales;
    \item $c(s,a,s') = d(s,s')$ es el costo de arista, igual a la distancia eucl\'idea
          entre celdas consecutivas ($1.0$ para movimientos ortogonales).
\end{itemize}

La funci\'on de transici\'on $\text{resultado}(s,a)=s+\Delta(a)$ es v\'alida si el destino
permanece dentro del tablero y no es un obst\'aculo. El espacio es finito, determinista,
totalmente observable y est\'atico.

\subsection{Dise\~no de la funci\'on heur\'istica}
\label{sec:heuristica}

Para A* se emple\'o la distancia eucl\'idea al objetivo:
\begin{equation}
\label{eq:heuristica}
h(n) = \sqrt{(x_n - x_g)^2 + (y_n - y_g)^2},
\end{equation}
con $n=(x_n,y_n)$ y el objetivo $g=(x_g,y_g)$. Esta heur\'istica es
\textbf{admisible}: en una cuadr\'icula con 4-conectividad cada paso cuesta al menos
$1$, de modo que el costo real restante es al menos la distancia Manhattan, y como
$h(n) \le d_{\text{Manhattan}}(n,g)$, se tiene $h(n) \le h^*(n)$ para todo $n$; nunca
sobreestima. Es adem\'as \textbf{consistente}, pues para toda arista $(n,n')$ de costo
$c(n,n')=1$,
\begin{equation}
\label{eq:consistencia}
h(n) \le 1 + h(n') \quad\Longleftrightarrow\quad h(n)-h(n') \le d(n,n')=1,
\end{equation}
desigualdad que se cumple por la desigualdad triangular de la distancia eucl\'idea. La
admisibilidad garantiza la \textbf{optimalidad} de A* y la consistencia evita
reexpansiones.

Para Hill Climbing se us\'o como funci\'on de evaluaci\'on el negativo de la misma
heur\'istica,
\begin{equation}
\label{eq:evaluacion}
f(n) = -\,h(n),
\end{equation}
de modo que maximizar $f(n)$ equivale a minimizar la distancia al objetivo. El algoritmo
solo eval\'ua los vecinos inmediatos del estado actual (\emph{steepest ascent}).

\subsection{Algoritmos implementados}
\label{sec:algoritmos}

A* se implement\'o con una cola de prioridad (\code{heapq}) ordenada por
$f(n)=g(n)+h(n)$, con conjuntos de abiertos y cerrados y un diccionario de predecesores
para reconstruir la ruta. Hill Climbing se implement\'o en su variante
\emph{steepest ascent} y, para mitigar los \'optimos locales, se le a\~nadi\'o
\textbf{reinicios aleatorios}: al estancarse, reinicia desde una celda libre aleatoria
conservando la mejor trayectoria. Ambos fueron implementados desde cero, usando
exclusivamente \code{numpy}, \code{pandas}, \code{matplotlib} y \code{heapq}. El
c\'odigo completo se incluye en el Anexo~\ref{anexo:codigo} y en el notebook
\code{T2Carlos\_Pi\~nero.ipynb}.

\section{Resultados Experimentales}
\label{sec:resultados}

La Tabla~\ref{tab:metricas} resume las m\'etricas. A* se ejecut\'o una vez (es
determinista); Hill Climbing se ejecut\'o $30$ veces con semillas distintas, por ser
estoc\'astico. Ambos resolvieron el problema con el mismo costo \'optimo de $24.00$
($25$ nodos), pero con un n\'umero de expansiones muy diferente.

\begin{table}[htbp]
\centering
\caption{Comparaci\'on de m\'etricas de desempe\~no sobre el mismo dataset.}
\label{tab:metricas}
\small
\begin{tabular}{lcc}
\hline
\textbf{M\'etrica} & \textbf{A*} & \textbf{Hill Climbing} \\
\hline
¿Lleg\'o al objetivo?            & S\'i        & S\'i ($100\%$ de $30$ corridas) \\
Longitud del camino (nodos)      & 25          & $25.0 \pm 0.0$ \\
Costo total (distancia eucl\'idea) & 24.00     & $24.00 \pm 0.00$ \\
Nodos expandidos                 & 152         & $59.0 \pm 0.0$ \\
Tiempo de ejecuci\'on (s)        & 0.001713    & 0.000575 \\
Reinicios aleatorios             & ---         & 0.0 \\
\hline
\end{tabular}
\end{table}

La Figura~\ref{fig:espacio} muestra el espacio de estados con los obst\'aculos y los
puntos de inicio y objetivo. La Figura~\ref{fig:astar} presenta la ruta \'optima de A*
junto con el orden de expansi\'on de los nodos. La Figura~\ref{fig:comparacion} compara
la ruta de A* con la mejor ruta hallada por Hill Climbing: n\'otese que, aunque son
distintas, ambas tienen costo $24.00$.

\begin{figure}[htbp]
\centering
\includegraphics[width=0.60\textwidth]{mapa_espacio_estados.png}
\caption{Espacio de estados cargado desde \code{dataset\_coordenadas.csv}.}
\label{fig:espacio}
\end{figure}

\begin{figure}[htbp]
\centering
\includegraphics[width=0.98\textwidth]{ruta_a_estrella.png}
\caption{A*: ruta \'optima (izquierda) y orden de expansi\'on de los $152$ nodos (derecha).}
\label{fig:astar}
\end{figure}

\begin{figure}[htbp]
\centering
\includegraphics[width=0.74\textwidth]{comparacion_rutas.png}
\caption{Superposici\'on de la ruta de A* y la mejor ruta de Hill Climbing.}
\label{fig:comparacion}
\end{figure}

La Figura~\ref{fig:analisis} analiza el costo computacional y la naturaleza estoc\'astica
de la b\'usqueda local. La Figura~\ref{fig:trampa} presenta el escenario adversario
descrito en la Secci\'on~\ref{sec:discusion}.

\begin{figure}[htbp]
\centering
\includegraphics[width=\textwidth]{analisis_optimos_locales.png}
\caption{Nodos expandidos, costo del camino y reinicios aleatorios de Hill Climbing.}
\label{fig:analisis}
\end{figure}

\section{An\'alisis y Discusi\'on}
\label{sec:discusion}

\subsection{\'Optimos locales, mesetas y la topolog\'ia del espacio}

Contrario a la intuici\'on inicial, en \emph{este} dataset Hill Climbing no fall\'o:
alcanz\'o el \'optimo global en las $30$ corridas y expandi\'o menos nodos que A*
($59$ frente a $152$). Para explicarlo es necesario analizar la topolog\'ia. El estado
objetivo se encuentra en la esquina inferior derecha, $(13,13)$, y el muro de
obst\'aculos situado en $x=5$ es \textbf{convexo respecto al objetivo}: en todo el
espacio transitable, moverse hacia abajo o hacia la derecha no aleja al agente de la
meta, y los obst\'aculos \emph{nunca} lo obligan a retroceder. En consecuencia, la
funci\'on de evaluaci\'on $f(n)=-h(n)$ decrece de forma mon\'otona y la trayectoria
voraz coincide con una ruta \'optima.

Para verificar que esto es una propiedad del tablero y no una casualidad, se ejecut\'o
Hill Climbing voraz \textbf{sin reinicios desde las $206$ celdas libres}: en el
$100\%$ de los casos se alcanz\'o el objetivo. Es decir, el paisaje de evaluaci\'on de
este problema \textbf{no posee m\'aximos locales, mesetas ni crestas} relevantes. Este
resultado refuerza una idea central: el desempe\~no de la b\'usqueda local depende por
completo de la forma del espacio, no de una propiedad intr\'inseca del algoritmo.

Para evidenciar la debilidad estructural de Hill Climbing se dise\~n\'o un
\textbf{escenario adversario controlado} (Figura~\ref{fig:trampa}): una concavidad en
forma de ``U'' que abre en direcci\'on \emph{opuesta} al objetivo. La evaluaci\'on
$f(n)=-h(n)$ conduce al agente hacia el fondo de la ``U'', donde todos los vecinos
empeoran; Hill Climbing se detiene tras solo dos pasos en el estado $(5,8)$, a una
distancia de $7.00$ de la meta: un \textbf{m\'aximo local}. A* bordea el obst\'aculo y
encuentra la soluci\'on con costo $14.00$. La regi\'on plana del fondo de la ``U''
corresponde adem\'as a una \textbf{meseta}: varios estados comparten el mismo valor de
evaluaci\'on, lo que en ausencia de movimientos laterales provoca el estancamiento.

\begin{figure}[htbp]
\centering
\includegraphics[width=0.68\textwidth]{escenario_trampa.png}
\caption{Escenario adversario en ``U'': HC queda atrapado en un \'optimo local mientras A* bordea el obst\'aculo.}
\label{fig:trampa}
\end{figure}

\subsection{Optimalidad y costo computacional de A*}

A* garantiz\'o la optimalidad \emph{por construcci\'on}: la heur\'istica de
la Ecuaci\'on~\eqref{eq:heuristica} es admisible (Ecuaci\'on~\eqref{eq:consistencia}),
por lo que A* expande el objetivo solo cuando su costo $g$ es \'optimo. Esta garant\'ia
es independiente de la topolog\'ia del tablero: incluso en el escenario adversario,
donde Hill Climbing fracasa, A* hall\'o la ruta \'optima.

El precio de la garant\'ia es computacional. A* expandi\'o $152$ nodos frente a los $59$
de Hill Climbing, un factor de aproximadamente $2.6\times$. La diferencia se debe a que
A* mantiene una \textbf{frontera} (cola de prioridad) que puede crecer de forma
exponencial, mientras que Hill Climbing solo conserva el estado actual y su vecindario.
En t\'erminos de complejidad, A* es $O(b^d)$ en tiempo y espacio en el peor caso, frente
al car\'acter de paso constante de la b\'usqueda local. Asimismo, A* requiere que
$h(n)$ sea admisible; una heur\'istica mal dise\~nada destruir\'ia la optimalidad.

Cabe se\~nalar que en este tablero la heur\'istica eucl\'idea es poco discriminante
(hay muchos empates de distancia), lo que obliga a A* a expandir m\'as nodos de los que
expandir\'ia una heur\'istica m\'as informada. Esto explica por qu\'e, aun siendo
\'optimo, A* result\'o m\'as ``costoso'' que la b\'usqueda local en este caso concreto.

\subsection{\'Cuando preferir b\'usqueda local sobre A*}

Hill Climbing resulta preferible cuando el espacio es enorme, la funci\'on de evaluaci\'on
es suave (sin \'optimos locales significativos) y basta una soluci\'on
\emph{suficientemente buena}. Ejemplos t\'ipicos incluyen la optimizaci\'on de
par\'ametros continuos, el ajuste de hiperpar\'ametros, la asignaci\'on de recursos en
tiempo real y la planificaci\'on con presupuestos de c\'omputo estrictos. En cambio, A*
se justifica cuando la \textbf{optimalidad} o la \textbf{completitud} son requisitos
cr\'iticos: navegaci\'on de robots, planificaci\'on de rutas, videojuegos con garant\'ias
de camino m\'inimo y, en general, cualquier escenario donde fallar o encontrar una ruta
sub\'optima tenga consecuencias.

\section{Conclusiones}
\label{sec:conclusiones}

La comparaci\'on emp\'irica permiti\'o constatar que A* es \'optimo y completo gracias a
la admisibilidad y consistencia de la distancia eucl\'idea, a costa de expandir m\'as
nodos y mantener una frontera en memoria. Hill Climbing, en cambio, result\'o
sorprendentemente eficiente en este tablero (mismo costo \'optimo con menos de la mitad
de expansiones), pero esa eficiencia dependi\'o de una topolog\'ia benigna: el barrido
desde las $206$ celdas libres confirm\'o la ausencia de \'optimos locales. El escenario
adversario en ``U'' demostr\'o la fragilidad estructural del m\'etodo, que qued\'o
atrapado en un m\'aximo local donde A* triunf\'o sin dificultad.

La lecci\'on principal es que la idoneidad de un algoritmo de b\'usqueda no puede
juzgarse en aislamiento, sino en funci\'on de la \textbf{topolog\'ia del espacio} y de
los \textbf{requisitos de garant\'ia}: A* cuando la optimalidad es obligatoria; b\'usqueda
local cuando la eficiencia prima y el paisaje de evaluaci\'on es favorable. Los reinicios
aleatorios mejoran la robustez de Hill Climbing pero incrementan el n\'umero de
evaluaciones y no restauran ninguna garant\'ia te\'orica.

\begin{thebibliography}{9}
\bibitem{russell2020} Stuart J. Russell and Peter Norvig.
\textit{Artificial Intelligence: A Modern Approach}. 4th Edition. Pearson, 2020.
\bibitem{hart1968} P. E. Hart, N. J. Nilsson, and B. Raphael.
\textit{A Formal Basis for the Heuristic Determination of Minimum Cost Paths}.
IEEE Transactions on Systems Science and Cybernetics, 4(2):100--107, 1968.
\end{thebibliography}

\appendix
\section{C\'odigo principal}
\label{anexo:codigo}

\begin{verbatim}
def distancia_euclidiana(a, b):
    return float(np.hypot(a[0] - b[0], a[1] - b[1]))

MOVIMIENTOS = [(1, 0), (-1, 0), (0, 1), (0, -1)]

def obtener_vecinos(nodo, datos):
    vecinos = []
    for dx, dy in MOVIMIENTOS:
        candidato = (nodo[0] + dx, nodo[1] + dy)
        if candidato in datos["libres"] or candidato == datos["objetivo"]:
            vecinos.append(candidato)
    return vecinos

def costo_camino(camino):
    return sum(distancia_euclidiana(camino[i], camino[i + 1])
               for i in range(len(camino) - 1))
\end{verbatim}

\begin{verbatim}
def a_estrella(datos):
    inicio, objetivo = datos["inicio"], datos["objetivo"]
    contador = itertools.count()
    abiertos = [(distancia_euclidiana(inicio, objetivo), next(contador), inicio)]
    g, padres, cerrados = {inicio: 0.0}, {inicio: None}, set()
    while abiertos:
        _, _, actual = heapq.heappop(abiertos)
        if actual in cerrados:
            continue
        cerrados.add(actual)
        if actual == objetivo:
            break
        for vecino in obtener_vecinos(actual, datos):
            if vecino in cerrados:
                continue
            nuevo_g = g[actual] + distancia_euclidiana(actual, vecino)
            if vecino not in g or nuevo_g < g[vecino]:
                g[vecino] = nuevo_g
                padres[vecino] = actual
                f = nuevo_g + distancia_euclidiana(vecino, objetivo)
                heapq.heappush(abiertos, (f, next(contador), vecino))
    # ... reconstruccion del camino mediante 'padres' ...
\end{verbatim}

\begin{verbatim}
def hill_climbing(datos, max_reinicios=50, max_iter=500, semilla=0):
    rng = random.Random(semilla)
    inicio, objetivo = datos["inicio"], datos["objetivo"]
    libres = list(datos["libres"])
    for reintento in range(max_reinicios):
        actual = inicio if reintento == 0 else rng.choice(libres)
        visitados = {actual}
        for _ in range(max_iter):
            if actual == objetivo:
                return True, visitados
            vecinos = [v for v in obtener_vecinos(actual, datos)
                       if v not in visitados]
            if not vecinos:
                break
            siguiente = min(vecinos,
                            key=lambda v: distancia_euclidiana(v, objetivo))
            if distancia_euclidiana(siguiente, objetivo) < \
                    distancia_euclidiana(actual, objetivo):
                actual = siguiente
                visitados.add(actual)
            else:
                break  # optimo local / meseta
    return False, visitados
\end{verbatim}

\end{document}
